% % % % % % % % % % % % % % % % % % % % % % % % % % %
% IS&T Template 
% Patrick Vandewalle
% January 2006
% % % % % % % % % % % % % % % % % % % % % % % % % % %

%%%%%%%%%%%%%%%%%%%%%%%%%%%%%%%%%%
% Document class
%%%%%%%%%%%%%%%%%%%%%%%%%%%%%%%%%%
\documentclass[letterpaper,twocolumn,fleqn]{article} 

%%%%%%%%%%%%%%%%%%%%%%%%%%%%%%%%%%
% Packages
%%%%%%%%%%%%%%%%%%%%%%%%%%%%%%%%%%
\usepackage{ist}
\usepackage[margin=1in]{geometry}
\usepackage{float}
\usepackage{wrapfig}
\usepackage{amsmath}
\usepackage{amssymb}
\usepackage{graphicx}
\usepackage{listings}
\usepackage{xcolor}
\usepackage{hyperref}
\usepackage{booktabs}
\pagestyle{empty}
\lstset{
  basicstyle=\ttfamily\footnotesize,
  breaklines=true,
  breakatwhitespace=true,
  columns=fullflexible,
  keepspaces=true,
  frame=single,
  showstringspaces=false,
  xleftmargin=0.02\columnwidth,
  xrightmargin=0.02\columnwidth
}

%%%%%%%%%%%%%%%%%%%%%%%%%%%%%%%%%%
% Title and Authors
%%%%%%%%%%%%%%%%%%%%%%%%%%%%%%%%%%

\title{Numerical Differentiation \& Integration Applied to\\
  Optical Flow and Movement Tracking in 2D Space}
\author{ Jake Davila, Jack Gaon, Jacky Li, Virgil Nhieu, Elijah Perez, Richie Wong\\ Cal Poly Pomona; Pomona, California \\ March 15, 2026}

%%%%%%%%%%%%%%%%%%%%%%%%%%%%%%%%%%
% Begin document
%%%%%%%%%%%%%%%%%%%%%%%%%%%%%%%%%%
\begin{document} 
\maketitle 
\thispagestyle{empty}

%%%%%%%%%%%%%%%%%%%%%%%%%%%%%%%%%%
% Abstract
%%%%%%%%%%%%%%%%%%%%%%%%%%%%%%%%%%

\begin{abstract}
Numerical differentiation and integration are foundational tools in
computational engineering.  In this report, we implement four forms of
single-variable quadrature (Trapezoidal, Simpson's~1/3, Simpson's~3/8,
Midpoint) and two forms of double integration (Trapezoidal~2D, Simpson~2D).
We apply these methods to an \textbf{optical flow} problem---tracking 2D
cursor movement along a figure-8 trajectory.  Velocity is estimated via
numerical differentiation of position data, and position is reconstructed by
integrating the estimated velocity.  A Lucas-Kanade-style matrix calculation
$(\mathbf{A}^T\mathbf{A})^{-1}\mathbf{A}^T\mathbf{b}$ demonstrates the
linear algebra behind real optical flow solvers.  Simpson's rules achieve
fourth-order convergence, dramatically outperforming the second-order
Trapezoidal and Midpoint rules as the number of steps increases from 10
to~1000.
\end{abstract}


\section{Introduction}

Numerical differentiation and numerical integration are fundamental tools in scientific computing. Numerical differentiation is used to approximate derivatives from discrete data, while numerical integration, also called quadrature, is used to approximate definite integrals. These methods are widely used in applications such as signal processing, control systems, and computer vision.

One important application in computer vision is \textbf{optical flow}, which describes the apparent motion of brightness patterns between consecutive image frames. The Lucas-Kanade method estimates this motion by combining three main numerical ideas:
\begin{enumerate}
  \item \textbf{Numerical differentiation} --- computing the spatial and temporal image gradients $I_x$, $I_y$, and $I_t$.
  \item \textbf{Matrix computation} --- solving the linear system $(\mathbf{A}^T\mathbf{A})\,\mathbf{d} = \mathbf{A}^T\mathbf{b}$ to estimate the local velocity vector.
  \item \textbf{Numerical integration} --- accumulating velocity over time to recover displacement and reconstruct the trajectory.
\end{enumerate}

In this report, four single-integral quadrature methods and two double-integral quadrature methods are implemented and compared in terms of convergence and accuracy. These methods are then applied to a 2D motion-tracking example to demonstrate how numerical differentiation, matrix-based estimation, and numerical integration work together in practice.


\section{Background}

\subsection{Four Single-Integral Quadrature Rules}

These numerical integration methods approximate the value of a definite integral by replacing the continuous area under a curve with weighted sums of function values at selected points. Each rule differs in accuracy and in how the sample points are used.

\subsubsection{Trapezoidal Rule}
The trapezoidal rule approximates the area under the curve by dividing the interval into small subintervals and treating each section as a trapezoid. It is simple and effective for smoothly varying functions.
\begin{equation}
  \int_a^b f(x)\,dx \approx \frac{h}{2}\bigl[f(x_0) + 2\textstyle\sum_{i=1}^{n-1} f(x_i) + f(x_n)\bigr]
  \label{eq:trap}
\end{equation}
Error: $O(h^2)$.

\subsubsection{Simpson's 1/3 Rule}
Simpson's 1/3 rule uses parabolic arcs instead of straight lines to approximate the function over pairs of subintervals. Because of this, it is generally more accurate than the trapezoidal rule for smooth data.
\begin{equation}
  \int_a^b f\,dx \approx \frac{h}{3}\bigl[f_0 + 4\textstyle\sum_{\text{odd}} f_i + 2\textstyle\sum_{\text{even}} f_i + f_n\bigr]
  \label{eq:s13}
\end{equation}
Error: $O(h^4)$. Requires $n$ even.

\subsubsection{Simpson's 3/8 Rule}
Simpson's 3/8 rule is another higher-order method based on cubic interpolation across three subintervals at a time. It is useful when the number of subintervals is a multiple of three.
\begin{equation}
  \int_a^b f\,dx \approx \frac{3h}{8}\bigl[f_0 + 3\textstyle\sum_{i\not\equiv 0} f_i + 2\textstyle\sum_{i\equiv 0} f_i + f_n\bigr]
  \label{eq:s38}
\end{equation}
Error: $O(h^4)$. Requires $n$ divisible by~3.

\subsubsection{Midpoint Rule}
The midpoint rule estimates the integral by evaluating the function at the center of each subinterval. It is often more accurate than using left or right endpoint estimates and works well for many smooth functions.
\begin{equation}
  \int_a^b f\,dx \approx h \sum_{i=0}^{n-1} f\!\left(\frac{x_i + x_{i+1}}{2}\right)
  \label{eq:mid}
\end{equation}
Error: $O(h^2)$.

\subsection{Double Integration}
To approximate a double integral over a rectangular region, the one-dimensional quadrature rules can be extended by applying them in both the $x$- and $y$-directions. This process is called tensor-product quadrature and is useful for estimating areas, volumes, or accumulated quantities over 2D domains.
\begin{equation}
  \iint_R f(x,y)\,dx\,dy \approx \sum_i \sum_j w_{ij}\, f(x_i, y_j)
  \label{eq:double}
\end{equation}

\subsection{Numerical Differentiation}
Numerical differentiation estimates the derivative of a function using nearby sampled values. This is especially useful when an exact derivative is difficult to obtain or when working with discrete experimental or simulated data.

Central finite difference:
\begin{equation}
  f'(x_i) \approx \frac{f(x_{i+1}) - f(x_{i-1})}{2h}
  \label{eq:central}
\end{equation}
Error: $O(h^2)$.

\subsection{Lucas-Kanade Optical Flow}
The optical flow constraint for a pixel is
\begin{equation}
  I_x u + I_y v + I_t = 0
  \label{eq:ofc}
\end{equation}
Over a window of $N$ pixels this yields the overdetermined system
$\mathbf{A}\,\mathbf{d} = \mathbf{b}$, where
$\mathbf{A} = [I_x\ I_y]$, $\mathbf{d} = [u\ v]^T$, $\mathbf{b} = -I_t$.
The least-squares solution is
\begin{equation}
  \mathbf{d} = (\mathbf{A}^T \mathbf{A})^{-1}\, \mathbf{A}^T \mathbf{b}
  \label{eq:lk}
\end{equation}


\section{Problem: Optical Flow on a Figure-8 Path}

We simulate a cursor tracing a figure-8:
\begin{align}
  x(t) &= \sin(t), \quad y(t) = \tfrac{1}{2}\sin(2t) \\
  v_x(t) &= \cos(t), \quad v_y(t) = \cos(2t)
\end{align}
over $t \in [0, 2\pi]$ with $N = 200$ samples.

The workflow is:
\begin{enumerate}
  \item \textbf{Differentiate} position to estimate velocity.
  \item \textbf{Integrate} velocity to reconstruct position.
  \item \textbf{Solve} the Lucas-Kanade matrix equation for a
        synthetic pixel window.
\end{enumerate}


\section{Implementation}

This section presents the numerical methods implemented in Python for integration, differentiation, and optical flow estimation. Each method is designed to support the motion-tracking analysis performed in this project.

\subsection{Single Integral (Trapezoidal)}
The trapezoidal integration function approximates the value of a definite integral by dividing the interval into \(n\) subintervals, evaluating the function at equally spaced points, and summing the areas of the resulting trapezoids. This provides a simple and efficient way to numerically estimate the accumulated quantity under a curve.

\begin{lstlisting}
def integrate_trapezoidal(self, f, a, b, n):
    h = (b - a) / n
    x = np.linspace(a, b, n + 1)
    y = np.array([f(xi) for xi in x])
    return h * (y[0]/2 + np.sum(y[1:-1]) + y[-1]/2)
\end{lstlisting}

\subsection{Numerical Differentiation}
The differentiation function computes the numerical derivative of discrete data using the central difference method for interior points. At the endpoints, forward and backward differences are used because neighboring points on both sides are not available. This method is used to estimate velocity components from sampled position data.

\begin{lstlisting}
def differentiate(self, y, h, method='central'):
    n = len(y); dy = np.zeros(n)
    for i in range(1, n - 1):
        dy[i] = (y[i+1] - y[i-1]) / (2*h)
    dy[0]  = (y[1] - y[0]) / h
    dy[-1] = (y[-1] - y[-2]) / h
    return dy
\end{lstlisting}

\subsection{Lucas-Kanade Matrix Solution}
The Lucas-Kanade implementation estimates motion by forming a linear system from the spatial image gradients \(I_x\), \(I_y\), and the temporal gradient \(I_t\). These values are assembled into a matrix equation, and the flow vector is obtained by solving the least-squares normal equations. This produces an estimate of the local optical flow in the horizontal and vertical directions.

\begin{lstlisting}
A = np.column_stack([Ix, Iy])  # Nx2
b = -It                        # Nx1
ATA = A.T @ A                  # 2x2
ATb = A.T @ b                  # 2x1
flow = np.linalg.inv(ATA) @ ATb
\end{lstlisting}


\section{Results}

\subsection{Single Integral Convergence}
Figure~\ref{fig:conv} shows the absolute error of all four methods as the
number of steps increases from 10 to 1000.  The Simpson rules (1/3 and 3/8)
converge as $O(n^{-4})$, while Trapezoidal and Midpoint converge as
$O(n^{-2})$.

\begin{figure}[H]
  \centering
  \includegraphics[width=0.48\textwidth]{unit03_single_integral_comparison.png}
  \caption{Error vs.\ number of steps for 4 quadrature methods (log-log).}
  \label{fig:conv}
\end{figure}

\subsection{Optical Flow: Differentiation and Integration}
Figure~\ref{fig:flow} shows: (a)~the true and reconstructed figure-8 path,
(b)~true vs.\ numerically estimated velocity, and (c)~position reconstruction
error.

\begin{figure}[H]
  \centering
  \includegraphics[width=0.48\textwidth]{unit03_optical_flow.png}
  \caption{Optical flow results: trajectory, velocity, and error.}
  \label{fig:flow}
\end{figure}

\subsection{Lucas-Kanade Matrix Calculation}
For a $5 \times 5$ synthetic pixel window with true velocity
$(u, v) = (1.5, -0.8)$:
\[
  \mathbf{A}^T\mathbf{A} =
  \begin{bmatrix} a_{11} & a_{12} \\ a_{21} & a_{22} \end{bmatrix},
  \quad
  \mathbf{d} = (\mathbf{A}^T\mathbf{A})^{-1}\,\mathbf{A}^T\mathbf{b}
\]
The estimated flow closely matches the true velocity, confirming the
least-squares formulation.


\section{Discussion}

\subsection{Quadrature Comparison}
\begin{table}[H]
\centering
\small
\begin{tabular}{@{}lcc@{}}
\toprule
\textbf{Method} & \textbf{Order} & \textbf{Constraint} \\
\midrule
Trapezoidal & $O(h^2)$ & None \\
Simpson's 1/3 & $O(h^4)$ & $n$ even \\
Simpson's 3/8 & $O(h^4)$ & $n \bmod 3 = 0$ \\
Midpoint & $O(h^2)$ & None \\
\bottomrule
\end{tabular}
\caption{Comparison of single-integral methods.}
\label{tab:methods}
\end{table}

Simpson's rules achieve two extra orders of accuracy at
essentially the same computational cost.  The trade-off is a constraint on~$n$.

\subsection{Optical Flow Pipeline}
The differentiate-then-integrate pipeline introduces small cumulative errors
(drift), typical in inertial navigation and visual odometry.  Higher-order
differentiation and integration methods reduce this drift.


\section{Conclusion}

We implemented quadrature and numerical differentiation methods for a 2D optical flow problem. Simpson's rules converge as $O(h^4)$, while Trapezoidal and Midpoint converge as $O(h^2)$. The Lucas-Kanade method recovers the velocity vector, and differentiation with integration forms a motion-tracking pipeline for robotics and computer vision.

\small
\begin{thebibliography}{9}
\bibitem{bib1}R. Burden, J. Faires, \emph{Numerical Analysis}, 10th ed.,
  Cengage, 2016.
\bibitem{bib2}B. Lucas, T. Kanade, ``An iterative image registration
  technique with an application to stereo vision,'' \emph{Proc.\ IJCAI},
  pp.\ 674--679, 1981.
\bibitem{bib3}B. Horn, B. Schunck, ``Determining optical flow,''
  \emph{Artificial Intelligence}, 17, 185--203, 1981.
\end{thebibliography}

\end{document}
